% Единственный источник содержания темы. UTF-8.
% Метки label и ключи в [...] сохраняют опубликованные ссылки.

\section{Теория}

\subsection{Сложение вероятностей}

Мы пока что рассматриваем только случай \textbf{конечного} \emph{пространства элементарных исходов}

\begin{theorem}[addition-of-probabilities/theorem/1]\label{theorem-1}
\textbf{Формула сложения вероятностей.} Для событий $A$ и $B$ верно:

\[\mathbb{P}(A\cup B) = \mathbb{P}(A) + \mathbb{P}(B) - \mathbb{P}(A\cap B)\]
\end{theorem}

\begin{proof}[addition-of-probabilities/theorem/1]\label{solution-addition-of-probabilities-theorem-1}
Разобьём событие $A$ на несовместные части $A\setminus B$ и $A\cap B$:

\[\mathbb P(A)=\mathbb P(A\setminus B)+\mathbb P(A\cap B).\]

Объединение $A\cup B$ состоит из несовместных событий $A\setminus B$ и $B$, поэтому

\[\mathbb P(A\cup B)=\mathbb P(A\setminus B)+\mathbb P(B).\]

Выразив из первого равенства $\mathbb P(A\setminus B)$ и подставив во второе, получаем

\[\mathbb P(A\cup B)=\mathbb P(A)+\mathbb P(B)-\mathbb P(A\cap B).\]

В обоих разбиениях использована аксиома аддитивности для несовместных событий.
\end{proof}

\begin{example}[addition-of-probabilities/theory/example]\label{task-example}
\textbf{Пример (объединение и пересечение событий).} Продавщица выставляет два костюма на две витрины. В магазине есть чёрные, синие и белые костюмы. Костюмов каждого цвета достаточно; на обе витрины можно выставить костюмы одного цвета. Событие $A$: белый костюм будет вывешен на первую витрину. Событие $B$: белый костюм будет вывешен на вторую витрину. Перечислите элементарные события, которые входят в $A \cap B,\ A\cup B$.
\end{example}

\begin{solution}[addition-of-probabilities/theory/example]\label{solution-addition-of-probabilities-theory-example}
Обозначим цвета буквами Ч, С, Б. Исход записываем двумя буквами: первая — цвет костюма на первой витрине, вторая — на второй. Один цвет может повторяться.

Тогда $A=\{\text{БЧ},\text{БС},\text{ББ}\}$ и $B=\{\text{ЧБ},\text{СБ},\text{ББ}\}$.

Оба события происходят одновременно только в исходе ББ:

\[A\cap B=\{\text{ББ}\}.\]

Хотя бы одно из них происходит в пяти исходах:

\[A\cup B=\{\text{БЧ},\text{БС},\text{ББ},\text{ЧБ},\text{СБ}\}.\]

Вероятности цветов в условии не заданы, поэтому здесь перечисляем события, не приписывая им численных вероятностей.
\end{solution}

\begin{definition}
\textbf{Несовместными} событиями называются такие события $A$ и $B$, что $A\cap B = \emptyset$
\end{definition}

\begin{example}[addition-of-probabilities/theory/exercise]\label{task-exercise}
\textbf{Упражнение.} Докажите, что вероятность невозможного события $(\emptyset)$ равна $0$.
\end{example}

\begin{solution}[addition-of-probabilities/theory/exercise]\label{solution-addition-of-probabilities-theory-exercise}
События $\Omega$ и $\varnothing$ несовместны, а их объединение равно $\Omega$. По аксиоме аддитивности

\[\mathbb P(\Omega)=\mathbb P(\Omega)+\mathbb P(\varnothing).\]

Вычитая $\mathbb P(\Omega)$ из обеих частей, получаем $\mathbb P(\varnothing)=0$.
\end{solution}

\begin{theorem}[addition-of-probabilities/theorem/2]\label{theorem-2}
\textbf{Формула сложения вероятностей несовместных событий.} Для \textbf{несовместных} событий $A$ и $B$ верно:

\[\mathbb{P}(A\cup B) = \mathbb{P}(A) + \mathbb{P}(B)\]
\end{theorem}

\begin{proof}[addition-of-probabilities/theorem/2]\label{solution-addition-of-probabilities-theorem-2}
Для несовместных событий $A\cap B=\varnothing$, а вероятность пустого события равна нулю. Поэтому из доказанной общей формулы сложения следует

\[\mathbb P(A\cup B)=\mathbb P(A)+\mathbb P(B)-0.\]

Полученное равенство также непосредственно совпадает с аксиомой аддитивности.
\end{proof}

\subsection{Благоприятствующие события}

\begin{definition}
Элементарные исходы, содержащиеся в событии $A$ называются \textbf{благоприятствующими событиями событию} $A$
\end{definition}

Теперь мы, зная вероятности элементарных исходов можем вычислять вероятности любых событий. Очевидно, что разные элементарные события несовместны.

\begin{theorem}[addition-of-probabilities/theorem/3]\label{theorem-3}
\textbf{Вероятность события} $A$  —  сумма вероятностей благоприятствующих событию $A$ элементарных исходов
\end{theorem}

\begin{proof}[addition-of-probabilities/theorem/3]\label{solution-addition-of-probabilities-theorem-3}
Пусть событию $A$ благоприятствуют элементарные события $\omega_1,\ldots,\omega_m$. Тогда $A$ — их объединение, а сами элементарные события попарно несовместны.

Повторное применение формулы сложения для несовместных событий даёт

\[\mathbb P(A)=\mathbb P(\omega_1)+\cdots+\mathbb P(\omega_m).\]

Если благоприятствующих событий нет, то $A=\varnothing$, и обе части равны нулю.
\end{proof}

\section{Задачи}

Первая часть учебника

\begin{task}[addition-of-probabilities/problems/1]\label{task-problems-1}
\textbf{1. №273} На тарелке лежат одинаковые на вид пирожки: 4 с мясом, 8 с капустой и 3 с вишней. Петя наугад берёт один пирожок. Найдите вероятность того, что пирожок окажется с вишней.
\end{task}

\begin{solution}[addition-of-probabilities/problems/1]\label{solution-addition-of-probabilities-problems-1}
Всего $4+8+3=15$ пирожков. Благоприятствуют три исхода — выбор каждого из трёх пирожков с вишней. Вероятность выбора каждого пирожка равна $\frac1{15}$, поэтому искомая вероятность равна $3\cdot\frac1{15}=\frac15=0,2$.
\end{solution}

\begin{task}[addition-of-probabilities/problems/2]\label{task-problems-2}
\textbf{2. №275} У бабушки десять чашек: шесть с красными цветами, остальные с синими. Бабушка наливает чай в случайно выбранную чашку. Найдите вероятность того, что это будет чашка с синими цветами.
\end{task}

\begin{solution}[addition-of-probabilities/problems/2]\label{solution-addition-of-probabilities-problems-2}
Синих чашек $10-6=4$. Благоприятствующие исходы — выбор любой из этих четырёх чашек. Из десяти равновозможных выборов подходят четыре:

\[\mathbb P=4\cdot\frac1{10}=\frac4{10}=0,4.\]
\end{solution}

\begin{task}[addition-of-probabilities/problems/3]\label{task-problems-3}
\textbf{3. №280} Симметричную монету бросают два раза. Найдите вероятность следующих событий:\\
 a) решка выпадет хотя бы один раз\\
 б) в первый раз выпадет орёл
\end{task}

\begin{solution}[addition-of-probabilities/problems/3]\label{solution-addition-of-probabilities-problems-3}
Четыре равновозможных исхода: ОО, ОР, РО, РР, где О — орёл, Р — решка. Вероятность каждого равна $\frac14$.

а) Хотя бы одна решка есть в исходах ОР, РО и РР. Вероятность равна $\frac34$.

б) Первый результат — орёл в исходах ОО и ОР. Вероятность равна $\frac24=\frac12$.
\end{solution}

\begin{task}[addition-of-probabilities/problems/4]\label{task-problems-4}
\textbf{4. №284} В некотором опыте возможно три элементарных события: $a, b$ и $c$. Вероятность того, что наступит либо событие $a$, либо событие $b$, равна $0,4$; вероятность того, что наступит либо событие $a$, либо событие $c$, равна $0,7$. Найдите вероятность каждого из элементарных событий.
\end{task}

\begin{solution}[addition-of-probabilities/problems/4]\label{solution-addition-of-probabilities-problems-4}
\figure{euler-284.svg}{Круги Эйлера: A = \{a, b\}, B = \{a, c\}; в пересечении исход a. Вероятности a, b, c равны 0,1; 0,3; 0,6.}{720}{430}

Положим $x=\mathbb P(a)$, $y=\mathbb P(b)$, $z=\mathbb P(c)$. Разные элементарные события несовместны, поэтому

\[x+y=0,4,\qquad x+z=0,7,\qquad x+y+z=1.\]

Вычитая первое равенство из третьего, получаем $z=0,6$. Вычитая второе из третьего, получаем $y=0,3$. Тогда $x=1-0,6-0,3=0,1$.

\textbf{Ответ:} $\mathbb P(a)=0,1$, $\mathbb P(b)=0,3$, $\mathbb P(c)=0,6$.
\end{solution}

\section{Домашнее задание}

Первая часть учебника, стр. 143–149. В каждой задаче надо кратко перечислить все благоприятствующие элементарные события.

\begin{task}[addition-of-probabilities/homework/1]\label{task-homework-1}
\textbf{1. №272} В шахматной партии Андрей играет с Борисом. Вероятность выигрыша Андрея равна 0,3, вероятность ничьей равна 0,2, вероятность того, что партия не будет закончена, равна 0,1. Найдите вероятность того, что:\\
 а) Андрей не проиграет;\\
 б) Борис не проиграет;\\
 в) никто не выиграет.
\end{task}

\begin{solution}[addition-of-probabilities/homework/1]\label{solution-addition-of-probabilities-homework-1}
Обозначим элементарные исходы: $a$ — выиграл Андрей, $b$ — выиграл Борис, $d$ — ничья, $u$ — партия не закончена. Тогда

\[\mathbb P(b)=1-0,3-0,2-0,1=0,4.\]

а) Андрей не проиграл в исходах $a,d,u$: $0,3+0,2+0,1=0,6$.

б) Борис не проиграл в исходах $b,d,u$: $0,4+0,2+0,1=0,7$.

в) Никто не выиграл в исходах $d,u$: $0,2+0,1=0,3$.

Здесь незавершённая партия считается отдельным исходом: в нём ни победа, ни поражение ещё не наступили.
\end{solution}

\begin{task}[addition-of-probabilities/homework/2]\label{task-homework-2}
\textbf{2. №277} В лыжных гонках участвуют 13 спортсменов из России, 2 спортсмена из Норвегии и 5 спортсменов из Швеции. Порядок, в котором спортсмены стартуют, определяется жребием. Найдите вероятность того, что первым будет стартовать спортсмен из России.
\end{task}

\begin{solution}[addition-of-probabilities/homework/2]\label{solution-addition-of-probabilities-homework-2}
Всего $13+2+5=20$ спортсменов. Элементарный исход — имя спортсмена, стартующего первым. Обозначим российских спортсменов $R_1,\ldots,R_{13}$: именно эти 13 исходов благоприятствуют событию.

Каждый спортсмен с вероятностью $\frac1{20}$ получает первый номер, поэтому

\[\mathbb P=13\cdot\frac1{20}=\frac{13}{20}=0,65.\]
\end{solution}

\begin{task}[addition-of-probabilities/homework/3]\label{task-homework-3}
\textbf{3. №282} Стрелок один раз стреляет в круглую мишень (рис. 65). Вероятности попадания в зоны мишени показаны в таблице 51. Число очков, которое получает стрелок, равно номеру зоны.
\end{task}

\begin{tabledata}{Таблица 51. Вероятности попадания в зоны мишени}
Зона & 1 & 2 & 3 & 4 & 5 \\
Вероятность & 0,001 & 0,002 & 0,004 & 0,006 & 0,021 \\
Зона & 6 & 7 & 8 & 9 & 10 \\
Вероятность & 0,065 & 0,138 & 0,243 & 0,334 & 0,186 \\
\end{tabledata}

\figure{target-zones.png}{Рисунок 65. Мишень с десятью зонами: голубые зоны 1, 5 и 9; жёлтые зоны 4 и 8}{530}{575}

Найдите вероятность события:\\
 а) «стрелок получит меньше 5 очков»;\\
 б) «стрелок получит больше 7 очков»;\\
 в) «стрелок попадёт в жёлтую зону»;\\
 г) «стрелок не попадёт в голубую зону».

\begin{solution}[addition-of-probabilities/homework/3]\label{solution-addition-of-probabilities-homework-3}
Элементарные исходы — попадания в зоны с номерами от 1 до 10. Суммируем вероятности нужных зон из таблицы 51.

а) Меньше 5 очков: зоны $\{1,2,3,4\}$.

\[\mathbb P=0,001+0,002+0,004+0,006=0,013.\]

б) Больше 7 очков: зоны $\{8,9,10\}$.

\[\mathbb P=0,243+0,334+0,186=0,763.\]

в) Жёлтые зоны на рисунке — $\{4,8\}$.

\[\mathbb P=0,006+0,243=0,249.\]

г) Не голубые зоны: $\{2,3,4,6,7,8,10\}$. Сумма их вероятностей равна

\[0,002+0,004+0,006+0,065+0,138+0,243+0,186=0,644.\]
\end{solution}

\begin{task}[addition-of-probabilities/homework/4]\label{task-homework-4}
\textbf{4. №283} В некотором опыте возможно три элементарных события: $a$, $b$ и $c$. Вероятность того, что наступит либо событие $b$, либо событие $c$, равна 0,83. Найдите вероятность элементарного события $a$.
\end{task}

\begin{solution}[addition-of-probabilities/homework/4]\label{solution-addition-of-probabilities-homework-4}
Искомому событию благоприятствует единственный элементарный исход $a$. События $a,b,c$ образуют всё пространство исходов, поэтому

\[\mathbb P(a)=1-\mathbb P(\{b,c\})=1-0,83=0,17.\]
\end{solution}

\begin{task}[addition-of-probabilities/homework/5]\label{task-homework-5}
\textbf{5. №285} Бросают одну игральную кость. Найдите вероятность события:\\
 а) «выпадет чётное число очков»;\\
 б) «выпадет число очков, кратное 3»;\\
 в) «выпадет больше 3 очков»;\\
 г) «выпадет число очков, кратное 7».
\end{task}

\begin{solution}[addition-of-probabilities/homework/5]\label{solution-addition-of-probabilities-homework-5}
Элементарные исходы — числа $1,2,3,4,5,6$, каждый имеет вероятность $\frac16$.

а) Благоприятствующие исходы $\{2,4,6\}$: $\mathbb P=\frac36=\frac12$.

б) Благоприятствующие исходы $\{3,6\}$: $\mathbb P=\frac26=\frac13$.

в) Благоприятствующие исходы $\{4,5,6\}$: $\mathbb P=\frac36=\frac12$.

г) Среди чисел от 1 до 6 нет кратных 7. Благоприятствующих исходов нет, поэтому вероятность равна 0.
\end{solution}

